\documentclass[11pt]{article}
\usepackage[margin=1in]{geometry}
\usepackage{graphicx}
\usepackage{wrapfig}
\usepackage{booktabs}

\title{Wrapped Figure Stress Case}
\author{Test Corpus}
\date{}

\begin{document}
\maketitle

\begin{abstract}
This fixture isolates a text-wrapped raster figure in a compact academic paper.
The image is deliberately small so that placement, caption flow, and paragraph
wrapping can be inspected without relying on a large external asset.
\end{abstract}

\section{Motivation}
\label{sec:motivation}

Empirical manuscripts often use small diagrams beside narrative explanations,
especially when the figure clarifies a classification rule or measurement
sequence. The conversion path should preserve the local relationship between the
figure and the surrounding paragraph because the meaning of the text depends on
that spatial arrangement.

\begin{wrapfigure}{r}{0.34\textwidth}
\vspace{-0.75\baselineskip}
\centering
\includegraphics[width=0.31\textwidth]{figures/wrapped_signal.png}
\caption{A compact diagnostic signal}
\label{fig:wrapped-signal}
\vspace{-0.5\baselineskip}
\end{wrapfigure}

The right-wrapped figure in Figure~\ref{fig:wrapped-signal} represents a small
diagnostic signal used during record review. The prose around the image is long
enough to flow above, beside, and below the figure, which creates a useful test
for line wrapping in conversion. A faithful output should keep the caption close
to the image and should avoid treating the surrounding sentences as an unrelated
block of text. The content is synthetic, but the layout resembles the compact
figures often placed near case definitions in applied research reports.

The review protocol classifies observations using a simple threshold. Values
above the horizontal reference line trigger a second check, while values below
the line remain in the primary analytic sample. This sentence extends the wrapped
region so that the final lines rejoin the full text width after the figure ends.

\section{Summary Table}
\label{sec:summary}

\begin{table}[htbp]
\centering
\caption{Synthetic review categories}
\label{tab:categories}
\begin{tabular}{lcc}
\toprule
Category & Records & Share \\
\midrule
Primary analytic sample & 184 & 0.74 \\
Secondary review & 51 & 0.21 \\
Deferred classification & 12 & 0.05 \\
\bottomrule
\end{tabular}
\end{table}

Table~\ref{tab:categories} supplies a conventional floating table after the
wrapped figure. The combination helps distinguish failures in ordinary floats
from failures specific to the wrapfig environment.

\end{document}
